\chapter{Presentation and Local Reconstruction}\label{chap:8}

\section*{Overview of the Chapter}

The question of this chapter is \textbf{when the overall geometry and the
structure-preserving morphisms can be recovered from local descriptions}.
For the changes preserving structure treated in Chapter~\ref{chap:7}, we treat the problem of
constructing a single morphism from the descriptions of its individual components.

As an example, consider the case where the modernization of an intrabank transfer system is
divided among the account, payment, and accounting teams.
The descriptions of the teams must agree on the shared transactions, amounts, and
completion conditions.
For the correspondence before and after the modernization as well, the correspondence
of transactions, the actions of the operations, and the evaluation of the Laws are
aligned across the teams.
In this chapter we construct the overall model from such divided descriptions and explain the
conditions under which the changes on the parts are assembled into a single
structure-preserving morphism.

We impose three conditions on reconstruction:
that the local data distinguish different overall morphisms,
that an overall morphism can be constructed from compatible correspondence tables,
and that an object can be constructed from a local description of the structure.
In this chapter we formulate these as separation, assembly of morphisms, and assembly of
objects, respectively, and prove them.
For full geometries we give these constructions from the primitive data of types, operations,
Laws, local geometry, and realization.
The main theorem shows that the category of full geometries and all their
structure-preserving morphisms is equivalent to the category of local models defined from
primitive data and compatibility conditions.

In the gluing of Chapter~\ref{chap:2} we fixed a geometry and a cover of contexts and glued
the states chosen on each context after correcting them.
There ``local'' refers to the range of the context in which a state is read.
In this chapter ``local'' refers to the readout of the finitely many types, values, and
equations that describe a geometry and a morphism.
The covers and the contexts themselves are included in the data that are read, and we assemble
the geometries equipped with them together with the structure-preserving morphisms.

What is used in the general reconstruction is a compatible family over all finite fragments.
Reconstruction requires specifying readout items sufficient to determine the objects and
morphisms, together with their compatibility conditions.
After the main theorem we treat the application to design changes and the cases where, by the laws of the operations, a finite
description suffices.
For finite models as well, we distinguish the choice of the necessary information from the
computational cost of checking and of reconstruction.

\begin{xltabular}{\linewidth}{@{}LLL@{}}
\toprule
Question & Construction & Outcome \\
\midrule
\endhead
When can objects and morphisms be finitely presented? & Finite exception codes and morphisms between fixed endpoints & Conditions for finite presentation and its limits \\
\addlinespace[3pt]
When can finite fragments be glued into a single table or function? & Restrictions and graphs of typed functions & Gluing of tables, and the conditions for being a function \\
\addlinespace[3pt]
What are the conditions for recovering the whole from local data? & Separation, assembly of morphisms, assembly of objects & Equivalence of categories given by the readout functor \\
\addlinespace[3pt]
From what are the geometries and morphisms of AAT assembled? & Primitive data of operations, Laws, geometry, and realization & Reconstruction of objects and of all structure-preserving morphisms \\
\addlinespace[3pt]
Can a design change be constructed from divided descriptions? & Local descriptions and correspondences for an intrabank transfer & Construction of the model and of the change from divided descriptions \\
\addlinespace[3pt]
Can an overall morphism be recovered from a finite table? & Unique extension by the laws of the operations & Conditions for recovering an overall morphism from a finite description \\
\addlinespace[3pt]
How can one check in a finite model? & Choice of readout items, enumeration and decision of equality & Distinction between the information sufficient for reconstruction and the computational procedure \\
\bottomrule
\end{xltabular}

The diagnostic invariance of Chapter~\ref{chap:3} guarantees that the verdict of a chosen
diagnosis does not change when the reading is changed.
The reflection of Chapter~\ref{chap:7} is the property that compatibility of the original
change follows from compatibility in the normalized image.
The reconstruction of this chapter recovers the objects and morphisms themselves.
Each of these properties holds for specified objects, morphisms, and observed quantities.

\section{Recovery of objects and morphisms by finite presentation}\label{sec:8.1}

In Chapter~\ref{chap:5} we showed that, if the presentations of the endpoints may be chosen,
a base morphism can be presented from a finite source and a finite or cofinite extraction.
When the presentations of the endpoints are fixed, not every semantic morphism can be
presented as a morphism between them.
In this section we use finite exception codes to make this difference explicit.

\subsection*{Finite exception tables and the one-point compactification}

\begin{definition}[Finite exception code]\label{def:8.1}

For a set $D$, we take a code to be a pair of a default value $b\in\{0,1\}$ and a finite set
$E\subseteq D$. The evaluation is
\begin{equation*}
\mathrm{ev}(b,E)(x)=
\begin{cases}
1-b,&x\in E,\\
b,&x\notin E
\end{cases}
\tag{8.1}\label{eq:8.1}
\end{equation*}
The default value 0 represents the characteristic function of a finite set, and the default
value 1 that of a cofinite set.

The default value specifies the value of the evaluation outside the finite exception set.
To represent it by a single point, we consider $D^+=D\sqcup\{\infty\}$.
We take each point of $D$ to be isolated, and we take a set containing $\infty$ to be open
if its complement is a finite subset of $D$.
This is the one-point compactification of the discrete space $D$.
When $D$ is finite, $\infty$ is isolated as well.

\end{definition}

\begin{proposition}[Topological meaning of codes]\label{prop:8.2}

Giving $\{0,1\}$ the discrete topology, we have the following bijection.
\begin{equation*}
\begin{aligned}
\mathrm{Code}(D)&\ \cong\ C(D^+,\{0,1\}),\\
(b,E)&\longmapsto
\left[
x\longmapsto
\begin{cases}
\mathrm{ev}(b,E)(x),&x\in D,\\
b,&x=\infty.
\end{cases}
\right]
\end{aligned}
\tag{8.2}\label{eq:8.2}
\end{equation*}
If $D$ is infinite, the code is determined uniquely by its evaluation on $D$ alone.
If $D$ is finite, then for every function $a:D\to\{0,1\}$ there are two codes, one with
default value 0 and one with default value 1.

\end{proposition}

\begin{proof}

The function obtained from a code is constant on the neighborhood $D^+\setminus E$ of
$\infty$ and is therefore continuous.
Conversely, for a continuous function $a$, put $b=a(\infty)$.
By continuity $a^{-1}(\{b\})$ is a neighborhood of $\infty$, so
$E=\{x\in D\mid a(x)\ne b\}$ is finite. The two constructions are mutually inverse.

For infinite $D$, a point can be taken outside the union of two finite exception sets.
If the evaluations agree, the default values agree there, and the exception sets agree as
well.
For finite $D$, the set $E_b=\{x\mid a(x)\ne b\}$ is finite for each $b$, and these give the
two codes.

\end{proof}

Theorem~\ref{thm:5.15} of Chapter~\ref{chap:5} therefore gives the following topological
characterization. Finiteness of the sources at both ends, together with the extension of each
extraction predicate of the target to a continuous function on $D^+$, characterizes the
finite presentation of a base morphism including the endpoint isomorphisms.
Finiteness of a discrete source is also equivalent, topologically, to compactness.
Indeed, that the open cover consisting of the individual points has a finite subcover means
that the total number of points is finite.
This characterization permits meaning-preserving changes of the enumeration of the
endpoints and of the correspondence of Atoms.

\subsection*{Morphisms between fixed presentations}

Next we investigate which semantic morphisms can be presented between fixed presentations.
We take as objects the codes of Chapter~\ref{chap:5} with a finite source, a normalization
table, and a finite exception table.
A morphism of presentations records a finite table of sources and a permutation moving only
finitely many Atoms.
We require that the normalization and the designated point be preserved and that the
transported extraction code be equal to the code of the target.
We identify the descriptions of permutations that give the same semantic morphism, and we
write $D_0$ for the decoding functor of the category obtained by this quotient.

We fix two presentations $P,Q$ and write $D_0P,D_0Q$ for their meanings.
We take $b_P(s)$ to be the default value of the extraction code used after normalizing the
source $s$.

\begin{proposition}[Existence condition for morphisms between fixed presentations]\label{prop:8.3}

Let $s_f$ be the source map and $\sigma_f$ the permutation of Atoms of a semantic base
morphism $f:D_0P\to D_0Q$. Then
\begin{equation*}
\exists h:P\to Q,\ D_0(h)=f
\quad\Longleftrightarrow\quad
\left\{
\begin{aligned}
&\{a\in D\mid \sigma_f(a)\ne a\}\text{ is finite},\\
&b_Q(s_f(s))=b_P(s)\quad\text{for all }s.
\end{aligned}
\right.
\tag{8.3}\label{eq:8.3}
\end{equation*}

\end{proposition}

\begin{proof}

A presented permutation moves only finitely many Atoms.
Moreover, the transport of a code does not change the default value, so the left-hand side
implies the right-hand side.

Conversely, assume the right-hand side.
Enumerating the finite set moved by $\sigma_f$ and writing the permutation on it as a table,
this permutation extends to $\sigma_f$ itself.
For the source map we use $s_f$.
Preservation of the normalization and of the designated point is a condition on $f$.
By exactness of the extraction, the evaluations of the transported code and of the target
code are equal.
Since the default values are equal as well, \eqref{eq:8.1} shows that the exception sets are
equal, and we obtain the equality of codes.

\end{proof}

The semantic morphisms that escape this presentation are therefore only those that move
infinitely many Atoms and those whose default values after normalization do not agree at
corresponding sources.
The next example shows that the latter occurs even for a finite set of Atoms.

\begin{example}[Non-fullness between fixed presentations]\label{ex:8.4}

We take $D=\{a\}$, a single source, and the identity normalization.
If the extraction codes are taken to be $(0,\varnothing)$ and $(1,\{a\})$, then both
extractions are always false.
The semantic identity morphism exists, but since the default values differ, that identity
morphism cannot be presented between these two codes.

This example shows that recovery of the Hom-sets between fixed presentations does not follow
from agreement of the meanings of the objects.
The information about the default values that remains in a presentation constrains the
presentability of semantic morphisms.

\end{example}

\begin{corollary}[Normalization of presentations over finitely many Atoms]\label{cor:8.5}

If $D$ is finite, each extraction code can be replaced by
\begin{equation*}
R_{\mathrm{fin}}(b,E)
=
\bigl(0,\{a\in D\mid\mathrm{ev}(b,E)(a)=1\}\bigr)
\tag{8.4}\label{eq:8.4}
\end{equation*}
On the full subcategory of the objects for which all the default values of the extraction
codes used after normalization are 0, the functor $D_0$ is fully faithful.
Moreover $D_0R_{\mathrm{fin}}\cong D_0$, and the semantic objects and morphisms are
preserved.

\end{corollary}

\begin{proof}

Equation~\eqref{eq:8.4} does not change the evaluation.
A permutation of a finite set has finite support and the default values at both ends are 0,
so Proposition~\ref{prop:8.3} applies to every semantic morphism.
Faithfulness comes from having identified the descriptions that give the same semantic
morphism.
The construction $R_{\mathrm{fin}}$ also acts on morphisms, retaining the source map and the
Atom permutation.
After evaluation each component is an identity isomorphism, and naturality with respect to
morphisms holds as well.

\end{proof}

\section{Compatibility of finite fragments and gluing}\label{sec:8.2}

In the previous section we saw that, even when the meanings of objects agree, the morphisms
between fixed presentations need not be recoverable.
In what follows we read the individual components of morphisms in addition to the objects,
and investigate the conditions under which the overall objects and morphisms are constructed
from compatible local descriptions.
In the gluing of finite fragments we distinguish the agreement of values on the common part
from the conditions under which the glued table defines a function or a structure.
We define the former as compatibility of finite fragments.
The latter includes the existence and uniqueness of an output for each input, as well as
preservation of the operations.

\subsection*{Indices and restrictions of the readout}

\begin{definition}[Compatible family of finite fragments]\label{def:8.6}

We take $\Lambda$ to be the set of the items to be read and $V_q$ to be the type of the value
of an item $q$.
For a finite subset $S\subseteq\Lambda$, we take a finite fragment to be an element of
\begin{equation*}
T_S=\prod_{q\in S}V_q
\tag{8.5}\label{eq:8.5}
\end{equation*}
For $S\subseteq T$ we use the restriction map $r_{S,T}:T_T\to T_S$ that forgets components.
A family $(a_S)_S$ over all the finite subsets is said to be \textbf{compatible} if
\begin{equation*}
r_{S,T}(a_T)=a_S
\quad\text{whenever }S\subseteq T
\tag{8.6}\label{eq:8.6}
\end{equation*}
holds.

What is finite here is the number of items contained in a single fragment.
The types of the values may include Booleans, elements of a ring, finite polynomials,
references to declared sets, and so on.
A finite fragment therefore does not necessarily represent a bit string of finite length.

\end{definition}

\begin{lemma}[Gluing from single points]\label{lem:8.7}

Between the overall table and the compatible families of finite fragments there is the
following bijection.
\begin{equation*}
\prod_{q\in\Lambda}V_q
\ \cong\
\left\{(a_S)_S\ \middle|\ r_{S,T}(a_T)=a_S\right\}.
\tag{8.7}\label{eq:8.7}
\end{equation*}

\end{lemma}

\begin{proof}

From the overall table we take the restrictions to the finite subsets.
In the other direction, we define the value at $q$ to be the single-point fragment
$a_{\{q\}}(q)$.
By \eqref{eq:8.6}, each component of any finite fragment agrees with this value.
Restriction and gluing compose to the identity in both directions.

\end{proof}

\subsection*{Representation by graphs of typed functions}

\begin{definition}[Graph of a typed function]\label{def:8.8}

We describe a function between declared sets $A,B$ by a Boolean table $R(x,y)$ on the pairs
of an input and an output.
When the candidate pairs of types are also included in the indices, we set the value to 0 on
the pairs other than the selected $(A,B)$.
Over the valid types we require the following.
\begin{equation*}
\forall x\in A,\quad\exists!y\in B,\quad R(x,y)=1.
\tag{8.8}\label{eq:8.8}
\end{equation*}
That is, each row has exactly one output.
We take the readout of a function $f:A\to B$ to be
$R_f(x,y)=1\Longleftrightarrow f(x)=y$.

\end{definition}

\begin{lemma}[Readout and assembly of graphs of functions]\label{lem:8.9}

The functions $A\to B$ and the tables satisfying Definition~\ref{def:8.8} are in one-to-one
correspondence.
This correspondence preserves identities and composition.

\end{lemma}

\begin{proof}

We take $f_R(x)$ to be the unique output in \eqref{eq:8.8}.
By definition $R_{f_R}=R$ and $f_{R_f}=f$.
The table of the identity is the diagonal $x=y$.
Composition is given by
\begin{equation*}
(R_g\star R_f)(x,z)=1
\quad\Longleftrightarrow\quad
\exists y,\ R_f(x,y)=1\ \land\ R_g(y,z)=1
\tag{8.9}\label{eq:8.9}
\end{equation*}
which is equivalent to $g(f(x))=z$.
By the uniqueness in each row, \eqref{eq:8.8} holds after composition as well.

\end{proof}

\begin{example}[Compatibility and the conditions for a graph of a function]\label{ex:8.10}

When $A$ is nonempty, the finite fragments of the Boolean table all of whose values are 0 are
compatible.
No row, however, has an output, so \eqref{eq:8.8} is not satisfied.
Moreover, placing two different outputs on one input satisfies existence but loses
uniqueness.
Compatibility of finite fragments alone therefore does not guarantee the conditions for a
graph of a function.

\end{example}

\begin{definition}[Local conditions by primitive evaluation]\label{def:8.11}

A formula of primitive evaluation is an equation or an implication assembled from
finitely many components of types, values, and graphs of functions.
We take the support of a formula to be the finite set of the indices referred to in it.
We take a local condition to be such a formula quantified over the specified inputs.

For example, for the actions of operations $d:A\to B$ and $d':A'\to B'$, the condition that
correspondences $h_A:A\to A'$ and $h_B:B\to B'$ preserve the actions,
\begin{equation*}
h_B(d(a))=d'(h_A(a))
\tag{8.10}\label{eq:8.10}
\end{equation*}
becomes, once $a$ and the intermediate output are fixed, a finite formula evaluating four
applications of functions.
Requiring this condition for all $a$ yields preservation of the actions of the operations.

Likewise, \eqref{eq:8.8} takes an output $y$ for each $x$ and requires its value to be 1 and
the value at every other candidate $z$ to be 0.
The formula of the evaluation for each candidate is finite, but the condition expressing the
totality of a row requires quantification.
When a finite carrier is required, the existence of a finite list covering all the elements
is also included in the conditions.
That each formula of evaluation has finite support is distinguished from the whole condition being verifiable by finitely many checks.

\end{definition}

\section{Reconstruction principle by separation and assembly}\label{sec:8.3}

Before turning to individual structures, we state a sufficient condition for local
reconstruction.
Let $N:\mathcal{R}\to\mathcal{M}$ be a functor that reads the overall structures and
morphisms into local models.

\begin{theorem}[Local reconstruction principle]\label{thm:8.12}

If $N$ satisfies the following three conditions, we obtain an equivalence of categories
$\mathcal{R}\simeq\mathcal{M}$ whose forward functor is $N$.

\begin{enumerate}
\item \textbf{Separation.} For $f,g$ with the same endpoints, $N(f)=N(g)$ implies $f=g$.
\item \textbf{Assembly of morphisms.} From any $u:NX\to NY$ we can construct
$\mathrm{asm}(u):X\to Y$ with $N(\mathrm{asm}(u))=u$.
\item \textbf{Assembly of objects.} From any $m\in\mathcal{M}$ we can construct an object
$A(m)\in\mathcal{R}$ and an isomorphism
$\varepsilon_m:N(A(m))\xrightarrow{\sim}m$.
\end{enumerate}

Then for each Hom-set we have
\begin{equation*}
\mathrm{Hom}_{\mathcal{R}}(X,Y)
\ \underset{\mathrm{asm}}{\overset{N}{\rightleftarrows}}\
\mathrm{Hom}_{\mathcal{M}}(NX,NY),\qquad
N\mathrm{asm}=1,\quad
\mathrm{asm}N=1
\tag{8.11}\label{eq:8.11}
\end{equation*}
and to each local morphism there corresponds exactly one overall morphism.

\end{theorem}

\begin{proof}

For a morphism $f:X\to Y$, assembly of morphisms gives $N(\mathrm{asm}(Nf))=Nf$.
Using separation we obtain $\mathrm{asm}(Nf)=f$, and \eqref{eq:8.11} holds.
Hence $N$ is fully faithful.
Assembly of objects gives essential surjectivity, so we obtain an equivalence of categories.
This is the standard characterization of an equivalence of categories
(\cite[\href{https://stacks.math.columbia.edu/tag/02C3}{Lemma 4.2.19, Tag 02C3}]{Stacks}).

We construct the inverse functor as follows.
To an object we assign $A(m)$, and for $u:m\to n$ we put
\begin{equation*}
A(u)=
\mathrm{asm}
\bigl(\varepsilon_n^{-1}\,u\,\varepsilon_m\bigr)
\tag{8.12}\label{eq:8.12}
\end{equation*}
Preservation of identities and composition holds in the image of $N$, and by separation it
holds in $\mathcal{R}$ as well.
By \eqref{eq:8.12}, $\varepsilon$ is a natural isomorphism.
The other natural isomorphism is given by the unique assembly of
$\varepsilon_{NX}^{-1}:NX\to N(A(NX))$.
Naturality of this natural isomorphism and the triangle identity on the $\mathcal{R}$ side
are checked by applying $N$ to both sides.
Using \eqref{eq:8.12} and the inverse laws of assembly and reducing the composites of
$\varepsilon$ with its inverse, the morphisms being compared agree.
By separation we obtain the original equality. The triangle identity on the $\mathcal{M}$
side follows from the fact that the composite of $\varepsilon_{NX}$ with its inverse is the
identity morphism.

\end{proof}

In this principle, separation guarantees the uniqueness of morphisms.
The existence of the morphisms and the objects requires proofs that construct the respective
structures from the local conditions.
In what follows we define local objects and local morphisms from types and equations and
prove these three conditions.

\section{Local models of full geometries}\label{sec:8.4}

\subsection*{The category of realizations and the variants of morphisms}

\begin{construction}[Input for full geometries and the category of realizations]\label{cons:8.13}

We fix a set of Atoms and a variant of morphisms defined below, and we take $\Theta$ to be
this pair.
We write $\mathcal{R}_\Theta$ for the category whose objects are the full geometries of
Chapter~\ref{chap:1} and whose morphisms are all the structure-preserving morphisms of the
chosen variant.
The local model is defined from the primitive data that constitute these objects and
morphisms.

Changes of full geometries arise in two ways: comparing the presentations after
transport by an equality, and retaining the concrete correspondence maps between
the two presentations. We use the following two variants, one for each.

\end{construction}

\begin{definition}[The two variants of morphisms of full geometries]\label{def:8.14}

In the \textbf{representative variant} we use the morphisms of Definition~\ref{def:1.30}.
They have maps of the base and of the core, transport of covers and overlaps, and a
homomorphism of coefficient rings, and we require the raw system of the codomain to be equal
to the raw system of the domain with its coefficients and contexts transported.
The comparison maps for supports, axes, and observables are natural with respect to the
restrictions chosen there.

In the \textbf{explicit-map variant} the components of the base, the core, the covers, the
overlaps, and the coefficients are the same, and the equality of raw systems is replaced by
the following concrete correspondence.
On each context after transport, the raw coordinates, the local data, and the relation
indices are put in bijection.
This correspondence preserves the evaluation of polynomials using the map of coefficients,
the relations, and the restrictions.
The image of each context morphism is also specified, and restriction morphisms are required
to be sent to restriction morphisms.
Bijections are specified for the supports, axes, and observables of each context as well, and
the predicates of the readout and naturality along these context morphisms are required.
For the raw coordinates, for example, we specify which coordinate a given coordinate is sent
to, and apply that correspondence to the relations and to the polynomials of the
restrictions.

In either variant, the maps of the base and of architecture objects and the map of
coefficients need not be invertible.
The identity is the identity in each component, and composition is the composition of the
corresponding components.

\end{definition}

\subsection*{Primitive data and local compatibility conditions}

\begin{definition}[Local objects and local morphisms of full geometries]\label{def:8.15}

We fix $\Theta$.
The following table collects the data that are read in order to reconstruct the objects and
morphisms of full geometries.
Each datum is read as a value or as one point of the graph of a function.
The sets that serve as the domain and the codomain of a function are specified by the type
reference of the corresponding role.

\begin{xltabular}{\linewidth}{@{}LLL@{}}
\toprule
Layer read & Primitive data read from an object & Primitive data read from a morphism \\
\midrule
\endhead
Base and Atoms & Types of the sources, the designated point, the graph of the normalization, the extraction predicates & Graphs of the source map and of the bijection of Atoms \\
\addlinespace[3pt]
architecture & Composition relations of finite Atom families, object formation, types and endpoints of the operations and their actions on Atoms & Graphs of the object map and of the operation map \\
\addlinespace[3pt]
Laws and invariants & Indices, both sides, and evaluations of the Laws, detectors, types and values of the invariants, signatures & Correspondence of indices, preservation of evaluations and of actions, correspondence of signatures \\
\addlinespace[3pt]
Contexts and observables & Types and relations of the contexts, restrictions, ring operations and readouts of the observables & Correspondence of contexts, ring isomorphisms of the observables and preservation of restrictions \\
\addlinespace[3pt]
Selected geometry & Coverage requirements, overlaps, coefficient ring, raw coordinates, relations, polynomials of the restrictions & Preservation of covers and overlaps, map of coefficients, and, according to the variant, the equality of raw systems or the graph of the correspondence \\
\addlinespace[3pt]
Realization & Types of supports, axes, and observables, readout predicates, actions of the context morphisms & Comparisons of each realization component and naturality \\
\bottomrule
\end{xltabular}

For example, one item of the action of an operation $u:A\to B$ represents the truth value of
the proposition that its morphism of configurations sends an Atom $a$ to an Atom $b$.
One item of the multiplication of the coefficient ring represents the truth value of the
equation $a\cdot b=c$.
The restrictions of the raw system are written as finite polynomials, and the types of their
variables and coefficients are matched up.
For a local value we specify these primitive components, and the geometry or the morphism as
a whole is constructed by assembling them.

In the readout of a morphism we distinguish the items of the domain object, of the codomain
object, and of the morphism itself.
Their common index set is
\begin{equation*}
\Lambda_\Theta
=
\Lambda_\Theta^{\mathrm{src}}
\sqcup
\Lambda_\Theta^{\mathrm{tgt}}
\sqcup
\Lambda_\Theta^{\mathrm{hom}}.
\tag{8.13}\label{eq:8.13}
\end{equation*}
We assign a type of values to each item and define the finite fragments and restrictions of
Definition~\ref{def:8.6}.
In the case of an object, only the items of the object are used.
For the domain and codomain parts of a morphism we require agreement with the specified local
objects at the two ends.

We take a \textbf{local object} to be a compatible family of finite fragments of an object
satisfying the following conditions.

\begin{itemize}
\item The types declared for each role agree, and each graph of a function satisfies
Definition~\ref{def:8.8}.
\item The laws of each structure constituting a full geometry of Chapter~\ref{chap:1} hold as
formulas of primitive evaluation.
\item The family of Atoms selected in the core is finite.
\end{itemize}

We take a \textbf{local morphism} to be a compatible family of finite fragments of a morphism
satisfying the agreement of the types and of the two ends, the totality and uniqueness of the
graphs of functions, and the preservation conditions listed in the table.
For the components that must be bijections we specify the graph in the opposite direction and
the two inverse laws.
On a ring homomorphism we impose preservation of 0, 1, addition, and multiplication, and on a
Law we impose the correspondence of the indices and preservation of the evaluations of both
sides.
Naturality is a componentwise equation for each specified restriction or context morphism.
The conditions on the polynomials of the raw system are evaluated from the correspondence of
finitely many coefficients and variables.

For example, writing $e$ for the correspondence of Atoms, $\Phi$ for the correspondence of
operations, and $d(u)$ for the morphism of configurations of an operation $u$, the condition that operations and their actions on configurations be preserved reads
\begin{equation*}
e\bigl(d(u)_{\mathrm{At}}(a)\bigr)
=
d(\Phi(u))_{\mathrm{At}}\bigl(e(a)\bigr)
\tag{8.14}\label{eq:8.14}
\end{equation*}
Naturality for the restrictions and the observables requires that applying the restriction and
the component maps in either order give equal values.
Expanding each formula at a fixed input gives finite support.
Requiring the formulas at all inputs yields the preservation conditions for the structure as
a whole.

In the invariant condition of Definition~\ref{def:1.28}, the bijection of value ranges is only
required to exist and is not included among the components retained by a morphism.
To express such an existence condition, we attach to the local presentation auxiliary graphs
representing the maps in both directions.
On these too we impose the types, the inverse laws, preservation of evaluations, and
compatibility with the restrictions of finite fragments.
After constructing a compatible presentation, we identify presentations that differ only in
the choice of the auxiliary graphs.
This quotient leaves the components of a morphism that are to be retained, such as the source,
the operations, the contexts, the coefficients, and the raw system.

\end{definition}

\begin{lemma}[The local category and the readout functor]\label{lem:8.16}

The local objects and local morphisms of Definition~\ref{def:8.15} form a category
$\mathcal{M}_\Theta$.
The readout of the primitive data defines a functor
\begin{equation*}
N_\Theta:\mathcal{R}_\Theta\longrightarrow\mathcal{M}_\Theta
\tag{8.15}\label{eq:8.15}
\end{equation*}

\end{lemma}

\begin{proof}

The identity of a local morphism is defined by the diagonal graph, and composition
componentwise by \eqref{eq:8.9}.
Each graph is again total and unique.
The preservation equation for a composite follows from two preservation equations, and the
graphs in the opposite direction preserve the inverse laws by being composed in the reverse
order.
The implications required for the covers and the restrictions are preserved as well, by
chaining two implications.
In the formulas of the raw system we compose the map of coefficients and the substitution of
variables in turn.

The correspondences in both directions can also be composed for the auxiliary presentations
of the invariants.
Presentations with equal retained components give composites with equal retained components,
so composition is still defined after the identification.
The unit and associativity laws follow from the composition of graphs and from componentwise
composition.

If each component of a realization is read, the types and laws hold by the definition of that
realization.
Since the readout of graphs of functions preserves identities and composition, these readouts
form a functor.

\end{proof}

\section{The local reconstruction theorem for full geometries}\label{sec:8.5}

From the primitive data and the conditions of Definition~\ref{def:8.15} we construct full
geometries and their structure-preserving morphisms.
With this construction we prove the conditions of separation and assembly of
Theorem~\ref{thm:8.12} and obtain the equivalence between the category of full geometries and
the category of local models.

\subsection*{Assembly of full geometries}

\begin{lemma}[Assembly and separation for full geometries]\label{lem:8.17}

For both variants of full geometries, a full geometry can be assembled from a local object.
Moreover, between any full geometries $G,H$, the readout and the assembly relating the local
morphisms $N_\Theta G\to N_\Theta H$ to the morphisms of full geometries $G\to H$ are
mutually inverse.

\end{lemma}

\begin{proof}

We first construct the objects in the order of dependence of the types and then assemble the
morphisms between those objects.
Finally we check that these constructions and the readout of the primitive data are mutually
inverse.

We glue the finite fragments by Lemma~\ref{lem:8.7}.
We obtain the types of the sources, the Atoms, the architecture objects, the configurations,
and the operations, construct from each graph, by Lemma~\ref{lem:8.9}, the normalization and
the actions of the operations on Atoms, and extract the values of the extraction, of the
composition relations of finite Atom families, and of the object formation.
The designated point belongs to the sources by the conditions on the types.
The conditions on the normalization and the extraction, the finiteness of the selected family
of Atoms, the conditions on the object formation and the composition relations, and the
conditions on the endpoints of the operations and on the morphisms of configurations give the
structures of the base and of the architecture of Chapter~\ref{chap:1}.

Next we extract the invariants and the signatures and the indices of the Laws.
We then construct the contexts and their relations and the ring operations and restrictions
of the observables of each context, and check their respective axioms. Using these contexts
and rings, we construct both sides of the Laws and their evaluations.
Rereading the primitive conditions as equations about the values of the assembled functions,
the conditions for the Laws and for detection hold.
This determines the core.

Over this core we take the coverage requirements and the overlaps and construct the selected
site.
The type and the operations of the coefficient ring satisfy the ring axioms.
We assemble the raw coordinates, the local data, the relation indices, and the polynomials of
the restrictions of each context.
The primitive conditions include preservation of the relations and the equations for the
identities and the composition of the restrictions, so these form a raw system.
Finally we combine the supports, axes, and observables with their readouts and with the
actions of the context morphisms, and obtain a full geometry.
By this order, the types referred to by the later data are already determined by the data
constructed earlier.

For morphisms as well, we first assemble the graph of each retained component into a function.
For the components that are bijections, the two inverse laws give the inverse map.
The local conditions for the sources, the normalizations, the extractions, the operations, the
Laws, and the observables give the preservation equations for the assembled maps as they
stand.
The same holds for preservation of the covers and the overlaps and for the homomorphism of
coefficients.
In the representative variant we obtain the equality of transported raw systems.
In the explicit-map variant we obtain the correspondence of raw systems and naturality along
the actual context morphisms.
The auxiliary graphs of the invariants give the existence of the required bijections of
values.
Changing the choice of the auxiliary graphs does not change the retained components of the
assembled morphism.

The graph of each function in this construction agrees with the original table by
Lemma~\ref{lem:8.9}.
Agreement of the finite fragments follows from Lemma~\ref{lem:8.7} as well.
Conversely, assembling the readout of a morphism of full geometries, the value of each
function agrees with the value of the original morphism.
Equality of morphisms is determined by the equality of these retained components, so the
result of the assembly is equal to the original morphism.
For objects as well, each structure is recovered, and we obtain a canonical isomorphism
through the identification of the types.

\end{proof}

\subsection*{The main theorem}

\begin{theorem}[Local reconstruction of full geometries]\label{thm:8.18}

We fix the set of Atoms and the variant of morphisms of Construction~\ref{cons:8.13}.
The primitive readout functor $N_\Theta$ into the local model of Definition~\ref{def:8.15}
gives an equivalence of categories.
\begin{equation*}
\mathcal{R}_\Theta\simeq\mathcal{M}_\Theta.
\tag{8.16}\label{eq:8.16}
\end{equation*}
Its forward functor is $N_\Theta$ of \eqref{eq:8.15} itself.
For any full geometries $X,Y$, the maps
\begin{equation*}
\mathrm{Hom}_{\mathcal{R}_\Theta}(X,Y)
\ \underset{\mathrm{asm}_\Theta}
{\overset{N_\Theta}{\rightleftarrows}}\
\mathrm{Hom}_{\mathcal{M}_\Theta}(N_\Theta X,N_\Theta Y)
\tag{8.17}\label{eq:8.17}
\end{equation*}
are bijections, and the readout and the assembly compose to the identity in both directions.

\end{theorem}

\begin{proof}

By Lemma~\ref{lem:8.17}, for either variant of morphisms the primitive readout separates
morphisms and all local morphisms and local objects can be assembled.
The three conditions of Theorem~\ref{thm:8.12} are therefore satisfied, and we obtain an
equivalence of categories whose forward functor is $N_\Theta$.

\end{proof}

The Hom-sets of the main theorem consist of all the structure-preserving morphisms,
including the
noninvertible ones permitted by the chosen variant.
Their existence is derived from the primitive types and the preservation conditions by
Lemma~\ref{lem:8.17}.
Objects are recovered up to isomorphism, and the morphisms between fixed endpoints are
recovered uniquely.

\begin{corollary}[Recovery of evaluations, composition, and isomorphisms]\label{cor:8.19}

For $u:N_\Theta X\to N_\Theta Y$ and
$v:N_\Theta Y\to N_\Theta Z$ we have
\begin{equation*}
\begin{aligned}
\mathrm{asm}_\Theta(1)&=1,\\
\mathrm{asm}_\Theta(vu)
&=\mathrm{asm}_\Theta(v)\mathrm{asm}_\Theta(u).
\end{aligned}
\tag{8.18}\label{eq:8.18}
\end{equation*}
The evaluation of each primitive component agrees with the value of the original local table.
Moreover, $N_\Theta(f)$ is an isomorphism if and only if $f$ is an isomorphism.

\end{corollary}

\begin{proof}

Reading both sides of \eqref{eq:8.18}, they become equal by functoriality of the readout and
by \eqref{eq:8.17}.
By separation the original morphisms are equal as well.
The agreement of the evaluations follows from the definition of the output chosen in the
assembly of graphs.
If a local morphism has an inverse, assembling that inverse and using \eqref{eq:8.18} gives
the inverse of the overall morphism.
The converse direction is functoriality.

\end{proof}

\section{Application to design changes}\label{sec:8.6}

\subsection*{Modernization of an intrabank transfer system}

We make the intrabank transfer of the opening concrete as an example explaining the
reconstruction of full geometries.
The division of business responsibilities that distinguishes account management, payment
execution, and accounting is also found in the
\href{https://bian.org/wp-content/uploads/2024/12/BIAN-Service-Landscape-V9_0-Value-Chain-View.pdf}{BIAN Service Domain Landscape 9.0}.
In this example we treat a transfer in a single currency without fees, and we set the
following states and Laws.

Suppose that the old design kept the information about the transfer, the update of the
accounts, and the posting in a single transaction record.
In the new design these are separated into a transfer record, an account statement, and a
journal entry, linked by a shared transaction reference.
Each team describes the following structures and the correspondence from the old design.

\begin{xltabular}{\linewidth}{@{}LLL@{}}
\toprule
Team & Structures described & Matters made compatible with the other descriptions \\
\midrule
\endhead
Account & Accounts, balances, debit and credit operations & Correspondence between transactions and balance changes, sending and receiving accounts of a transfer, amount \\
\addlinespace[3pt]
Payment & Transfer requests, amounts, processing states, acceptance and completion operations & The transaction referred to by each record, the completion condition of a transfer \\
\addlinespace[3pt]
Accounting & Journal entries, transaction references, posting and correction operations & Correspondence between transactions and journal entries, agreement of amounts, balance of debits and credits \\
\bottomrule
\end{xltabular}

The division among the teams is a division of descriptions, and a finite fragment of this
chapter consists of the finitely many types, values, and equations that each team describes.
After declaring the shared types of transactions and amounts, we specify the endpoints and
actions of the operations, the evaluations of the Laws, the contexts and their restrictions,
and the data of the selected geometry and of the realization, following
Definition~\ref{def:8.15}.
When the same readout item appears in several descriptions, its values must agree.

We exhibit part of the compatibility conditions through a Law about the completion of a
transfer.
We fix a transfer identifier $q$ and a positive amount $m$ in a single currency.
We take $c,d,i,j\in\{0,1\}$ to be the values representing, respectively, the completion
verdict, the record of the debit from the sending account, the record of the credit to the receiving account, and the existence of the corresponding pair of debit and credit journal
entries.
For the last three, the determination also requires that the transaction reference and
the amount agree with $q,m$, respectively.
The completion condition of this example is $c=1\Longrightarrow d=i=j=1$.
Evaluating the values as integers, this can be expressed as the condition that the three
residuals $c(1-d)$, $c(1-i)$, and $c(1-j)$ are all 0.
Since each residual refers to finitely many values, it can be described as a primitive
evaluation in the sense of Definition~\ref{def:8.11}.

Suppose, for example, that for a completed transaction of the old design the payment team and
the accounting team specify the completed and posted records of the new design, while the
account team specifies the records before the debit and the credit.

\begin{xltabular}{\linewidth}{@{}LLLL@{}}
\toprule
Readout item & Transaction in the old design & Incompatible candidate correspondence & Candidate correspondence satisfying this Law \\
\midrule
\endhead
Completion verdict $c$ & 1 & 1 & 1 \\
\addlinespace[3pt]
Debit record $d$ & 1 & 0 & 1 \\
\addlinespace[3pt]
Credit record $i$ & 1 & 0 & 1 \\
\addlinespace[3pt]
Pair of journal entries $j$ & 1 & 1 & 1 \\
\bottomrule
\end{xltabular}

For the incompatible candidate, the first two residuals are 1.
Even though the types of the references to each record are correct, this candidate does
not preserve the evaluation of the Law required for a completed transfer.
The rightmost candidate satisfies this Law. To construct an overall morphism we further
require that the actions of the operations and all the preservation conditions listed in
Definition~\ref{def:8.15} be satisfied.

In the assembly of objects we glue these primitive data by Lemma~\ref{lem:8.7} and construct
the full geometry of the new design by Lemma~\ref{lem:8.17}.
Once the necessary data are collected and a local object is defined, its realization is
obtained together with an isomorphism to the readout.
In the assembly of morphisms we fix the set of Atoms and the variant of morphisms as a common
input $\Theta$.
We write $G_{\mathrm{old}},G_{\mathrm{new}}$ for the full geometries of the old and the new
design and specify the correspondences of the teams as a local morphism
$u:N_\Theta G_{\mathrm{old}}\to N_\Theta G_{\mathrm{new}}$.
Theorem~\ref{thm:8.18} gives a unique structure-preserving morphism
$h:G_{\mathrm{old}}\to G_{\mathrm{new}}$ whose readout is this morphism.
The objects and morphisms constructed here are those of an AAT model in which the data of
types, operations, Laws, local geometry, and realization are made explicit.

\begin{figure}[htbp]
\centering
% Figure 8.1 -- local reconstruction of an AAT model and a design change.
% Same content as the Japanese manuscript's figure ../../figures/ch08-local-reconstruction.svg, redrawn in TikZ.
\begin{tikzpicture}[
    x=1mm, y=1mm,
    team/.style={rounded corners=1.5mm, draw=black!45, semithick, fill=black!3,
                 align=center, font=\footnotesize, inner sep=2mm, text width=44mm,
                 minimum height=26mm,
                 execute at begin node={\hyphenpenalty=10000 \exhyphenpenalty=10000 }},
    wide/.style={rounded corners=1.5mm, draw=black!60, semithick, fill=white,
                 align=center, font=\footnotesize, inner sep=2mm, text width=136mm},
    result/.style={rounded corners=1.5mm, draw=black!45, semithick, fill=black!3,
                align=center, font=\footnotesize, inner sep=2mm, text width=62mm,
                minimum height=17mm},
    arr/.style={-{Stealth[length=5pt]}, semithick},
    head/.style={font=\small\bfseries},
    note/.style={font=\small}
  ]
  \node[head] at (75,97) {Local descriptions for an internal transfer};
  \node[note] at (75,90) {Declared structures and old-to-new correspondences};
  \node[team] (acct) at (25,72)
    {\textbf{Account management}\\[1pt]
     Accounts and balance changes\\ Debit / credit operations\\ Account correspondences};
  \node[team] (pay) at (75,72)
    {\textbf{Payment management}\\[1pt]
     Requests, amounts and states\\ Acceptance / completion\\ Transaction correspondences};
  \node[team] (jrnl) at (125,72)
    {\textbf{Accounting}\\[1pt]
     Journal entries and references\\ Posting / correction operations\\ Journal-entry correspondences};
  \node[wide] (cond) at (75,44)
    {\textbf{All required primitive data and compatibility conditions}\\[1pt]
     Shared readings agree on types, transaction references, amounts and completion\\
     Actions, laws, geometry and realizations satisfy the declared conditions};
  \node[result] (obj) at (39,17)
    {\textbf{Object assembly}\\[1pt] An AAT model, up to isomorphism};
  \node[result] (mor) at (111,17)
    {\textbf{Morphism assembly}\\[1pt]
     A unique structure-preserving map\\ between fixed old and new AAT models};
  \draw[arr] (acct.south) -- (acct.south |- cond.north);
  \draw[arr] (pay.south)  -- (pay.south  |- cond.north);
  \draw[arr] (jrnl.south) -- (jrnl.south |- cond.north);
  \draw[arr] (obj.north |- cond.south) -- (obj.north);
  \draw[arr] (mor.north |- cond.south) -- (mor.north);
  \node[note] at (75,3) {AAT local reconstruction: Theorem 8.18};
\end{tikzpicture}

\caption{\textbf{Descriptions of the primitive data and the correspondences by the account,
payment, and accounting teams.}
From all the required data and the compatibility conditions, objects are reconstructed up to
isomorphism and the morphisms between fixed endpoints are reconstructed uniquely.
In addition to the agreement of the types of the individual records, we require preservation of the actions of the operations, of the evaluations of the Laws,
and of the selected geometry and the realization.}\label{fig:8.1}
\end{figure}

By this construction, the compatibility of a design change divided among teams can be
described as the existence of an overall structure-preserving morphism.
For another modernization proposal as well, we can fix the same old design and the structures
to be retained and compare the proposals by checking whether the local correspondences to that
proposal satisfy the preservation conditions.

\subsection*{Application to normalization and generated comparisons}

\begin{example}[Tagged operations]\label{ex:8.20}

The operation $(u,t)$ of Example~\ref{ex:6.25} carries an action $u$ on configurations
together with a tag $t\in\{0,1\}$ that does not affect that action.
We take $\Omega$ to be the set of all architecture objects and put, for
$\chi:\Omega\to\{0,1\}$,
\begin{equation*}
T_\chi(u,t)=(u,t+\chi(A))
\quad\text{for }u:A\to B
\tag{8.19}\label{eq:8.19}
\end{equation*}
Addition is modulo 2.
Since the Laws, the invariants, and the geometry of this example do not depend on the value of
the tag, $T_\chi$ is an automorphism of the full geometry in the explicit-map variant.
The uniform tag flip is $\chi=1$.
We call $\chi$ a selection of flips at each starting object.

Theorem~\ref{thm:8.18} recovers all these $T_\chi$, together with the canonical normalization
of the same example, as morphisms of the original full geometry.
By reading the tag component, two morphisms with equal actions on configurations can also be
distinguished.
The equations for the composition with the normalization hold in the local category as they
stand as well, by Corollary~\ref{cor:8.19}.

\end{example}

\begin{example}[Objects, morphisms, and changes of endpoints for generated comparisons]\label{ex:8.21}

We use the common input of Construction~\ref{cons:5.37}. We fix the square realized by finite
codes and, for the same comparison diagram, the core, the strong edges, the conditions on the
end of the face and on the two routes in the base, and the specified comparisons.
On this basis we build a comparison from the same face, cochain, coefficients, and original
full geometry.

We take $G_{\mathrm d},G_{\mathrm v}$ to be the full geometries of the two routes and
$\overline\alpha:G_{\mathrm d}\xrightarrow{\sim}G_{\mathrm v}$ to be the canonical
comparison.
Writing $e,d$ for the normalizations of Chapter~\ref{chap:6} at the two ends, the generated
comparison satisfies
\begin{equation*}
\overline\beta=d\overline\alpha,\qquad
e^2=e,\quad d^2=d,\qquad
d\overline\alpha=\overline\alpha e
\tag{8.20}\label{eq:8.20}
\end{equation*}
Hence $d\overline\beta=\overline\beta=\overline\beta e$.

These are objects and morphisms of the representative variant, and Theorem~\ref{thm:8.18}
recovers all of the original geometry used in the generation, the geometries of the two
routes, and $\overline\alpha,\overline\beta,e,d$.
It also recovers all the automorphisms at the two ends and the commutative squares between
them.
By Corollary~\ref{cor:8.19}, \eqref{eq:8.20} is preserved in the local model as well.

\end{example}

\section{Application of the reconstruction principle to the semantics of CS}\label{sec:8.7}

We apply the local reconstruction principle of Theorem~\ref{thm:8.12} to the lenses and the
protocols with observations introduced in Chapter~\ref{chap:1}.
We first recover the objects and the general morphisms from the primitive data and the
preservation conditions.
Next, using the respective laws, we exhibit the reconstruction from a finite description, the
reference fiber or the table of generating edges.

\subsection*{Local models and assembly}

\begin{construction}[Local models for the semantics of CS]\label{cons:8.22}

We choose the input $\Theta$ from the following table and write $\mathcal{R}_\Theta$ for the
category of realizations.

\begin{xltabular}{\linewidth}{@{}LLL@{}}
\toprule
Input & What is fixed & Category of realizations $\mathcal{R}_\Theta$ \\
\midrule
\endhead
lens & The set $V$ of views and a reference value $v_0\in V$ & The lenses that satisfy the three laws and have a finite reference fiber, and all the maps preserving get and put \\
\addlinespace[3pt]
protocol & A finite presentation $Q$ and an observation functor $O$ on the quotient path category & The realizations whose state set at each vertex is finite, and all the natural transformations preserving the observations \\
\bottomrule
\end{xltabular}

The objects and morphisms are defined from the semantics of Definitions~\ref{def:1.32}
and~\ref{def:1.36}, respectively.
Their operations and Laws are connected to descriptions by Atoms and operations through the
typed construction of \S\ref{sec:1.10}.
An arbitrary map between reference fibers, and a noninjective state map, are also morphisms if
they satisfy the respective preservation conditions.

For the local description we use the following primitive data.

\begin{xltabular}{\linewidth}{@{}LLL@{}}
\toprule
Input & Data read from an object & Data read from a morphism \\
\midrule
\endhead
lens & The type of the states, the graphs of get and of each put & The graph of the state map \\
\addlinespace[3pt]
protocol & The type of the states at each vertex, the graphs of the generating edges and of the observations & The graphs of the maps at each vertex \\
\bottomrule
\end{xltabular}

Types are specified by the reference for their role, and we take each value of a graph of a
function to be a readout item.
Among the items of a morphism we distinguish the source, the target, and the morphism itself,
and we use the finite fragments and restrictions of \S\ref{sec:8.2}.
We take a local object to be a compatible family of finite fragments satisfying the agreement
of the types and the totality and uniqueness of each graph.
On lenses we impose the three laws and the finiteness of the reference fiber, and on protocols
the relations of paths, the naturality of the observations, and the finiteness of the state
set at each vertex, as conditions of primitive evaluation.
Finiteness is expressed by the existence of a finite list covering all the elements of the set
in question.

On a local morphism we require agreement with the specified objects at the two ends and the
totality and uniqueness of the graphs.
For lenses we impose preservation of get and put, and for protocols preservation of the
generating edges and of the observations, by componentwise equations.
The preservation equations for the operations and the observations have finite support once
the input is fixed.
The diagonal graph and the composition of \eqref{eq:8.9} preserve these conditions, so we
obtain a local category $\mathcal{M}_\Theta$.
We write $N_\Theta:\mathcal{R}_\Theta\to\mathcal{M}_\Theta$ for the functor that reads each
component of a realization.

\end{construction}

\begin{lemma}[Assembly and separation for lenses]\label{lem:8.23}

From the local objects and local morphisms for lenses, the lenses and their general morphisms
can be assembled.
The readout and the assembly are mutually inverse on each Hom-set, and the objects are
recovered up to isomorphism as well.

\end{lemma}

\begin{proof}

From a local object we take the state set $C$ and assemble the graph of get and the graph of
put at each view $v$.
We write $g:C\to V$ and $p:C\times V\to C$ for the resulting maps.
The three pointwise conditions of the local object are
\begin{equation*}
p(c,g(c))=c,\qquad
g(p(c,v))=v,\qquad
p(p(c,v),w)=p(c,w)
\tag{8.21}\label{eq:8.21}
\end{equation*}
Hence $(C,g,p)$ is a total lens.
The local object also carries the condition that a finite list covering $g^{-1}(v_0)$ exists,
so its reference fiber is finite.

We assemble the graph of a local morphism into $h:C\to C'$.
The two local conditions of a morphism give
\begin{equation*}
g'h=g,\qquad
h(p(c,v))=p'(h(c),v)
\tag{8.22}\label{eq:8.22}
\end{equation*}
so $h$ is a morphism of lenses.
By Lemmas~\ref{lem:8.7} and~\ref{lem:8.9}, the two directions are mutually inverse for
all the graphs and finite fragments.

\end{proof}

\begin{lemma}[Assembly and separation for protocols]\label{lem:8.24}

From the local objects and local morphisms for protocols with observations, the realizations
and their general morphisms can be assembled.
The readout and the assembly are mutually inverse on each Hom-set, and the objects are
recovered up to isomorphism as well.

\end{lemma}

\begin{proof}

For each vertex $v$ of the finite presentation $Q$ we obtain the state set $X_v$.
By the condition that a finite list covers all the states, $X_v$ is finite.
We assemble the graph of a generating edge $e:v\to w$ into $X_e:X_v\to X_w$ and the graph of
the observation into $o_v:X_v\to O(v)$.

A local object satisfies the equality of actions for each specified relation $p=q$.
Moreover, the local condition at each generating edge gives
\begin{equation*}
o_wX_e=O(e)o_v
\tag{8.23}\label{eq:8.23}
\end{equation*}
By Proposition~\ref{prop:1.37}, there is therefore a unique realization $(X,o)$ with these
finite states, edge maps, and observations.

From a local morphism we obtain the vertex maps $h_v:X_v\to Y_v$.
The local conditions are
\begin{equation*}
h_wX_e=Y_eh_v,\qquad
o_v^Yh_v=o_v^X
\tag{8.24}\label{eq:8.24}
\end{equation*}
where the first equation holds at all the generating edges and the second at all the vertices.
By Proposition~\ref{prop:1.38}, these vertex maps determine a unique semantics-preserving
morphism $h:X\Rightarrow Y$.
All the components are obtained uniquely from the graphs, and the extension is unique as
well.
Since the maps of Lemmas~\ref{lem:8.7} and~\ref{lem:8.9} are mutually inverse, so are
the readout and the assembly.

\end{proof}

\begin{corollary}[Local reconstruction for the semantics of CS]\label{cor:8.25}

For each input of Construction~\ref{cons:8.22}, the primitive readout functor $N_\Theta$ gives
an equivalence of categories $\mathcal{R}_\Theta\simeq\mathcal{M}_\Theta$.
The objects are recovered up to isomorphism, and between any fixed realizations the readout
and the assembly are mutually inverse, including for noninvertible morphisms.
This reconstruction preserves identities, composition, and the primitive evaluations, and
reflects isomorphisms.

\end{corollary}

\begin{proof}

Lemmas~\ref{lem:8.23} and~\ref{lem:8.24} each give the three conditions of
Theorem~\ref{thm:8.12}.
We therefore obtain an equivalence of categories whose forward functor is the primitive
readout.
Preservation of identities and composition and reflection of isomorphisms follow from
functoriality of the readout and from the mutual inverses on each Hom-set.
Agreement of the primitive evaluations follows from the mutual inverses of the readout and the
assembly in Lemmas~\ref{lem:8.7} and~\ref{lem:8.9}.

\end{proof}

\subsection*{Reconstruction of lenses from the reference fiber}

We restate Proposition~\ref{prop:1.34} of Chapter~\ref{chap:1} as the table used for the
reconstruction of morphisms and its extension.

\begin{proposition}[Tables on the reference fiber and general morphisms of lenses]\label{prop:8.26}

Let $L=(C,g,p)$ and $L'=(C',g',p')$ be lenses with the same views $V$ and the same reference
value $v_0$.
Put $K=g^{-1}(v_0)$ and $K'=(g')^{-1}(v_0)$.
Every map $t:K\to K'$ extends uniquely to a morphism of lenses by
\begin{equation*}
\mathrm{ext}(t)(c)
=
p'\bigl(t(p(c,v_0)),\,g(c)\bigr)
\tag{8.25}\label{eq:8.25}
\end{equation*}
Here $p(c,v_0)$ is read as an element of $K$.
This extension and the restriction to the reference fiber are mutually inverse.

\end{proposition}

\begin{proof}

Equation~\eqref{eq:8.25} is the extension \eqref{eq:1.19} of Proposition~\ref{prop:1.34}.
By the same proposition, this extension preserves get and put and is inverse to the
restriction to the reference fiber.

\end{proof}

The table needed for the reconstruction of a morphism is a map between $K,K'$, and finiteness
of the set $V$ of views is not assumed.
The correspondence at each view is determined as the unique extension by the law for put.

\begin{example}[Merging complementary information]\label{ex:8.27}

We take the state of a product lens to be $(a,k)\in V\times K$, take get to read the first
component, and take put to replace the first component by the specified view.
From a map $t:K\to K'$ of reference fibers we obtain
\begin{equation*}
h(a,k)=(a,t(k))
\tag{8.26}\label{eq:8.26}
\end{equation*}
For example, taking $K=\{0,1,2\}$ and $K'=\{0,1\}$ with $t(0)=t(1)=0$ and $t(2)=1$ gives a
noninjective morphism merging two values of the complementary information.
Equation~\eqref{eq:8.26} preserves get and put and, by Proposition~\ref{prop:8.26}, is
determined uniquely by this table on the reference fiber.

\end{example}

\subsection*{Reconstruction of protocols from the generating edges}

\begin{proposition}[Unique extension of an adapter from the generating edges]\label{prop:8.28}

For two realizations $(X,o^X),(Y,o^Y)$ of a protocol with observations, finite tables
$h_v:X_v\to Y_v$ at each vertex extend to a general morphism if and only if they satisfy the
first equation of \eqref{eq:8.24} at all the generating edges and the second at all the
vertices.
The extension is unique, and restriction and extension are mutually inverse.

\end{proposition}

\begin{proof}

Necessity follows from the naturality and the preservation of observations of
Definition~\ref{def:1.36}.
Sufficiency and uniqueness of the extension are obtained by applying
Proposition~\ref{prop:1.38} to the vertex tables $h_v$.
The vertex components of the resulting morphism are the original tables themselves, so
restriction and extension are mutually inverse.

\end{proof}

Preservation of observations is required for isolated vertices with no generating edges as
well.

For lenses, a map on the reference fiber determines, by the law for put, a morphism on all the
states.
For protocols, the maps at each vertex determine, by naturality at the generating edges, a
natural transformation on all the paths.
In both cases the laws of the operations are used to extend a local map to an overall
morphism.

\begin{example}[Preservation of observations and naturality]\label{ex:8.29}

Consider a protocol with two vertices and one generating edge $e:v\to w$.
We take the state set at each vertex of both realizations to be $\{0,1\}\times\{0,1\}$, and
the observation functor assigns $\{0,1\}$ to both vertices and the identity map to the edge.
We take the observation of each state to be $(a,k)\mapsto a$.
We define the actions of the edge on states by $X_e(a,k)=(a,k)$ and $Y_e(a,k)=(a,0)$.
Taking the identity map at both vertices as the candidate correspondence, the observations are
preserved.
At $k=1$, however, $h_wX_e(a,1)=(a,1)$ while $Y_eh_v(a,1)=(a,0)$, so naturality fails.
The preservation condition at the generating edges thus detects a difference of actions that
the observations alone cannot distinguish.

\end{example}

\subsection*{Compatibility of finite presentations with the primitive readout}

\begin{corollary}[Compatibility with the finite decoding functor]\label{cor:8.30}

For lenses and for protocols there are, respectively, the following category
$\mathcal{P}_\Theta$ of finite tables and decoding functor
$D_\Theta:\mathcal{P}_\Theta\to\mathcal{R}_\Theta$.

We take the presentation objects for lenses to be the natural numbers $n$ and the morphisms to
be arbitrary tables $\mathrm{Fin}(n)\to\mathrm{Fin}(m)$.
The decoding functor sends $n$ to the product lens $V\times\mathrm{Fin}(n)$ and a table $t$ to
$(v,k)\mapsto(v,t(k))$.

We take the presentation objects for protocols to be a natural number $n_v$ at each vertex, a
table for each generating edge, and a table of the observations, satisfying the relations and
the naturality of the observations.
A morphism is a vertex table preserving the generating edges and the observations.
The decoding functor extends the tables of the generating edges to the composites of paths.

Both decoding functors are equivalences of categories, and the composite with the primitive
readout of Corollary~\ref{cor:8.25},
\begin{equation*}
\mathcal{P}_\Theta
\xrightarrow{D_\Theta}
\mathcal{R}_\Theta
\xrightarrow{N_\Theta}
\mathcal{M}_\Theta
\tag{8.27}\label{eq:8.27}
\end{equation*}
is an equivalence of categories as well.
Moreover, the route that assembles the local model from the primitive data and then reads the
reference fiber or the table of generating edges agrees, up to natural isomorphism, with the
route that reads them directly from the realization.

\end{corollary}

\begin{proof}

For lenses, by Proposition~\ref{prop:8.26} the decoding functor is a bijection on each
Hom-set.
For objects, combining the isomorphism of Proposition~\ref{prop:1.33},
\begin{equation*}
C\xrightarrow{\sim}V\times K,\qquad
c\longmapsto\bigl(g(c),p(c,v_0)\bigr)
\tag{8.28}\label{eq:8.28}
\end{equation*}
with an enumeration of the finite set $K$, every lens becomes isomorphic to an image of the
decoding functor.
For protocols we enumerate the finite state set at each vertex and transport the generating
edges and the observations to that enumeration.
Proposition~\ref{prop:8.28} gives the bijection of the Hom-sets, so we obtain an equivalence
of categories in the same way.
Equation~\eqref{eq:8.27} is the composite of these with Corollary~\ref{cor:8.25}.

To exhibit the last natural isomorphism, we take $F_\Theta$ to be the functor that reads the
reference fiber or the table of generating edges and $A_\Theta$ to be the inverse functor of
the equivalence of Corollary~\ref{cor:8.25}.
Applying $F_\Theta$ to the unit isomorphism $1\cong A_\Theta N_\Theta$ of
Theorem~\ref{thm:8.12}, we obtain
\begin{equation*}
F_\Theta
\ \cong\
F_\Theta A_\Theta N_\Theta
\tag{8.29}\label{eq:8.29}
\end{equation*}
Hence the two routes agree not only on the objects but also on the general morphisms.

\end{proof}

The procedure for obtaining a morphism from a finite presentation differs between the two
semantics.
For lenses we take as input a table between the finite reference fibers and construct by
\eqref{eq:8.25} a map on all the states.
Preservation of get and put is guaranteed on all the states by Proposition~\ref{prop:8.26}.
When the set of views is infinite, the preservation conditions for an arbitrarily given map on all
the states need not be checkable on finitely many states alone.

For protocols we enumerate the finite states at each vertex and check, for a candidate vertex
table, the naturality of \eqref{eq:8.24} at all the generating edges and the preservation of
observations at all the vertices.
When the presentation object itself is also checked, in addition to the tables of the edges
and the naturality of the observations, we compare the actions of paths of finite length for
the specified finitely many relations.
If the evaluation of the necessary tables and the decision of equality are computable, these
can be carried out as finite checks.

\begin{proposition}[Connection with retracts, the Karoubi envelope, and comparison squares]\label{prop:8.31}

For both inputs, lenses and protocols, the category of realizations is idempotent complete.
Writing $D_\Theta$ for the finite decoding functor, \eqref{eq:8.27} extends to
\begin{equation*}
\mathrm{Kar}(\mathcal{P}_\Theta)
\simeq
\mathcal{R}_\Theta
\simeq
\mathcal{M}_\Theta
\tag{8.30}\label{eq:8.30}
\end{equation*}
and agrees, up to natural isomorphism, with the original decoding functor on the inclusion of
the finite presentations.
Moreover, these equivalences extend to the arrow categories and to the commutative squares of
comparisons, and they are compatible with the evaluation at the endpoints and with the
exchange of the Karoubi envelope and the arrow category of Chapter~\ref{chap:6}.

\end{proposition}

\begin{proof}

Idempotent completeness of lenses over a fixed view is given by Proposition~\ref{prop:6.30}.
For protocols, for an idempotent semantics-preserving adapter $p:X\Rightarrow X$ we take
\begin{equation*}
X_v^p=\{x\in X_v\mid p_v(x)=x\}
\tag{8.31}\label{eq:8.31}
\end{equation*}
By Proposition~\ref{prop:6.31}, this image, with the executions and observations restricted,
is a realization of the same protocol, and $p_v$ together with the inclusion gives its
splitting.

By Corollary~\ref{cor:8.30}, every realization is isomorphic to an image of the finite decoding
functor and is in particular a retract of that image.
Using the fully faithful decoding functor and the splitting above, sending $(P,e)$ to the split
image of $D_\Theta(e)$ yields \eqref{eq:8.30} from the Karoubi envelope of
Chapter~\ref{chap:6}.
When $e=1$, the result is isomorphic to the decoded original object.

The action on the arrow categories is defined by applying the equivalence functors at both ends
of a morphism and of a commutative square.
Compatibility with the evaluation at the endpoints follows from this definition.
The equivalence
$\mathrm{Kar}(\mathrm{Arr}(\mathcal{C}))\simeq\mathrm{Arr}(\mathrm{Kar}(\mathcal{C}))$ of
Chapter~\ref{chap:6} also uses the same idempotents at both ends and the same squares.
Taking the arrow category after the Karoubi envelope and taking the Karoubi envelope after
the arrow category give the same data of endpoints, morphisms, and squares.
The differences coming from the choice of the splittings agree under that canonical
isomorphism.

\end{proof}

\section{Description and reconstruction in finite models}\label{sec:8.8}

\subsection*{Resource constraints of implementations and finiteness of models}

In an implementation whose storage is fixed at $b$ bits, the states that can be stored in that
memory number at most $2^b$.
Integers of fixed bit width, and strings with an upper bound on the character set and on the
stored length, also form finite sets.
Under such resource constraints, the problem of reconstruction becomes the question of which
information a specified structure and change can be recovered from.

To use the finiteness of an implementation for reconstruction from finite tables, the types,
coefficients, and contexts adopted by the model and the range of the readout items must also
be specified.
In what follows we separate the case of tabulating all the required items from the case of
recovering by the laws from a partial readout.

\subsection*{The overall table and partial readouts}

If the index set $\Lambda$ in Definition~\ref{def:8.6} is finite, a compatible family consists
of the table $a_\Lambda$ containing all the items and its restrictions.
The gluing of Lemma~\ref{lem:8.7} is the recovery of this overall table.
When each value can be encoded finitely, the overall table also has a finite description.

Even when the items are finitely many, a chosen part alone need not distinguish the
candidates.
In the intrabank transfer example of \S\ref{sec:8.6}, if only the completion verdict $c$ and
the existence $j$ of journal entries are read, the incompatible candidate and the candidate
satisfying the Law both return $(c,j)=(1,1)$.
Adding the debit record $d$ and the credit record $i$, we distinguish $(d,i)=(0,0)$ for the
former from $(d,i)=(1,1)$ for the latter and can check the completion condition.
This shows that, even for finite candidates, the choice of the items read governs the outcome
of the check.

Theorem~\ref{thm:8.18} recovers the overall structure and morphisms using all the necessary
primitive data and preservation conditions.
To recover a morphism between fixed endpoints from fewer items, the values of the omitted
components must be determined uniquely by the laws.
The reconstructions from the reference fiber and from the table of generating edges in the
previous section are examples of such a unique extension.

\subsection*{Finite checking procedures and computational cost}

\begin{xltabular}{\linewidth}{@{}LLL@{}}
\toprule
Property & Content & Counterpart in this chapter \\
\midrule
\endhead
Reconstruction from all the required data & Recovers objects up to isomorphism and morphisms between fixed endpoints uniquely from descriptions satisfying the compatibility conditions & Theorem~\ref{thm:8.18} \\
\addlinespace[3pt]
Determination by a chosen finite table & The table extends uniquely to an overall morphism & Reconstruction from the reference fiber and from the table of generating edges \\
\addlinespace[3pt]
Finite executable procedure & Computes the check of the conditions and the extension using explicit enumerations and decisions of equality & The finite presentations of \S\ref{sec:8.7} \\
\bottomrule
\end{xltabular}

If the tables in question and the conditions to be checked can be enumerated finitely and each
evaluation and decision of equality can be computed, then whether the specified conditions hold
can be decided by finitely many checks.
The existence of a finite set alone does not amount to giving these enumerations or
computational procedures.

Moreover, a check terminating in finitely many steps is different from its being executable
within given time and memory resources.
A practical check requires examining the size of the input tables and the cost of evaluating
each condition.
The reconstruction theorems of this chapter show the existence and uniqueness of the
reconstruction and do not give upper bounds on the computational cost.

\section{Supplement: presentation and reconstruction when infinite objects are allowed}\label{sec:8.9}

As a mathematical generalization, we treat the case where infinite sets are allowed for the
Atoms and for the indices of the readout.
Specification models that abstract away resource bounds sometimes use such sets.
The two counterexamples below are statements about the morphisms and the readouts that such a
model permits.
They do not mean that presentation or reconstruction is impossible for an implementation under
fixed finite resources.

\subsection*{Limits of presentation by a countable syntax}

In the first counterexample we consider all the automorphisms of an infinite set and
impose no computability restriction.
Using the cardinality of that set, we show that a decoding map from a countable syntax is not
surjective.

\begin{proposition}[Impossibility of a surjective presentation by a countable syntax]\label{prop:8.32}

The base category $B$ of Chapter~\ref{chap:1} has an object whose set of automorphisms does not
fit into the image of the decoding map from any countable set of syntax.

\end{proposition}

\begin{proof}

We take the Atoms to be $\mathbb{N}$, the source to be a single point, the normalization to be
the identity, and the extraction to be always true.
We write $b_\infty$ for this base object.
For any $S\subseteq \mathbb{N}$ we define the permutation $\sigma_S$ that exchanges $2n$ and
$2n+1$ exactly when $n\in S$.
Since the extraction is constant, each $\sigma_S$ is an automorphism of $b_\infty$.
The set of exchanged pairs determines $S$ uniquely, so
\begin{equation*}
\mathcal P(\mathbb{N})\hookrightarrow\mathrm{Aut}_{B}(b_\infty),
\qquad S\longmapsto\sigma_S
\tag{8.32}\label{eq:8.32}
\end{equation*}
is injective. The left-hand side is uncountable, so there is no surjection from a countable set
onto the right-hand side.

\end{proof}

For example, the set of all finite strings over countably many symbols is countable, and this
conclusion applies to it.
This argument by cardinality depends on imposing no computability condition and on allowing all
subsets $S\subseteq \mathbb{N}$.
From this proposition one cannot conclude the impossibility of presentation when the scope is restricted to computable changes described by finite codes.

Moreover, the permutation $\sigma_{\mathbb{N}}$ exchanging all the adjacent pairs has infinite
support, so by Proposition~\ref{prop:8.3} it cannot be presented as a morphism of finite
exception codes.
On the other hand, this permutation can be described by the finite rule of adding 1 to even
numbers and subtracting 1 from odd numbers.
Presentability in a specified tabular format and describability by a finite rule are different.

This example holds for an object of the base category. The selected extraction is infinite, and the family of Atoms selected in the core is
therefore not finite.
In the local descriptions from \S\ref{sec:8.2} on, sets are retained as type references and the
whole index set is allowed to be infinite.
Reconstruction by this readout and surjective presentation by a countable syntax are different
properties.

\subsection*{Limits of finite fragments for infinite indices}

The second counterexample uses the tag change of Example~\ref{ex:8.20}.
For each $A$ in the set $\Omega$ of architecture objects, whether the tags of the operations
starting at $A$ are flipped is determined independently of the choices at the other objects.
We recover the map $\chi:\Omega\to\{0,1\}$ representing this choice from all the finite
fragments.

\begin{proposition}[Recovery of a tag change from all the finite fragments]\label{prop:8.33}

The selections of flips at each starting object $\chi:\Omega\to\{0,1\}$ of Example~\ref{ex:8.20} and the
compatible families of Boolean tables $a_S$ on the finite sets $S\subseteq\Omega$ are in
one-to-one correspondence.
With addition modulo 2, this becomes the group isomorphism
\begin{equation*}
\{0,1\}^{\Omega}
\ \cong\
\varprojlim_{S\subseteq\Omega,\ S\text{ finite}}\{0,1\}^{S}
\tag{8.33}\label{eq:8.33}
\end{equation*}
If $\Omega$ is infinite, the readout on any single finite set $S$ cannot distinguish all these
selections of flips at each starting object.

\end{proposition}

\begin{proof}

The bijection is Lemma~\ref{lem:8.7} applied to $V_q=\{0,1\}$.
Addition and restriction are defined pointwise, so the bijection preserves the group
operation.
Moreover, \eqref{eq:8.19} gives $T_\psi T_\chi=T_{\chi+\psi}$.
Since $\chi(A)$ is determined by the action on the operation obtained by attaching the tag 0
to the identity action of each $A$, this family is also represented faithfully as a subgroup of
the actual automorphisms.

If $\Omega$ is infinite and $S$ is finite, take $A\notin S$ and compare the function that is
identically 0 with the function that is 1 only at $A$.
Their values on $S$ are equal, but their actions on the tags of the operations starting at $A$
differ.

\end{proof}

If $\Omega$ is finite, this family can be determined by reading the whole of it with
$S=\Omega$.
The counterexample in the infinite case uses the fact that outside every finite fragment there
remains a place that can be changed.
Unlike the first counterexample, which uses cardinality, this argument needs only a change that
differs at a single point outside the range of the readout.

When applying these counterexamples, therefore, one must specify the infinite sets in the
model, the range of the changes allowed, and the readout scheme.
The questions of sufficiency of information and of computational cost for finite models are
treated by the conditions stated in \S\ref{sec:8.8}.

\section*{Summary of the Chapter}

In this chapter we gave a method for recovering objects and all the structure-preserving
morphisms between them from local descriptions of geometries and morphisms.

\begin{itemize}
\item \textbf{The scope of finite presentation.} With finite exception codes, even when the
meaning of an object can be recovered, the morphisms between fixed presentations need not
represent the whole semantic Hom-set.
Over finitely many Atoms, normalization of the presentations resolves this non-fullness
(\S\ref{sec:8.1}).
\item \textbf{Separation and assembly.} From typed readout items, compatibility of finite
fragments, and the totality and uniqueness of graphs of functions, we assembled tables and
functions.
Once separation of morphisms, assembly of morphisms, and assembly of objects are all in place,
the readout functor gives an equivalence of categories (\S\S\ref{sec:8.2}--\ref{sec:8.3}).
\item \textbf{Recovery of full geometries.} We proved the three conditions for the primitive
data of types, operations, Laws, local geometry, and realization.
The equivalence of categories of the main theorem recovers the objects and all the
structure-preserving morphisms of the chosen variant, including the noninvertible ones, and
preserves identities and composition (\S\S\ref{sec:8.4}--\ref{sec:8.5}).
\item \textbf{Finite descriptions by the laws of the operations.} The same reconstruction
principle applies to lenses and protocols as well.
From a finite reference fiber or a table of generating edges, a general morphism extends
uniquely by the laws of the operations (\S\ref{sec:8.7}).
\item \textbf{Sufficiency of information and computation.} For finite models we distinguish the
information that determines objects and morphisms, the procedure based on enumeration and
decision of equality, and the computational cost.
When infinite objects are allowed, a difference appears between reconstruction from all the
finite fragments and presentation by a single finite table or a countable syntax
(\S\S\ref{sec:8.8}--\ref{sec:8.9}).
\end{itemize}

In the modernization of an intrabank transfer system, once the descriptions of the account, payment,
and accounting teams define a local object that supplies the necessary primitive data
and satisfies the conditions, the whole design is constructed up to isomorphism.
Between a fixed old design and a fixed new design, when the correspondences of the teams
satisfy all the preservation conditions, they can be assembled into a unique
structure-preserving morphism.
Hence, if the same old design and the structures to be retained are fixed, several
modernization proposals can be compared from the point of view of whether a change that is
compatible overall exists.
